\begin{lemma}
\label{lemma-characteristic-polynomial-prime}
Let $R \to A$ be a ring homomorphism.
Assume $A \cong R^{\oplus n}$ as an $R$-module.
Let $f \in A$. The multiplication map $m_f: A
\to A$ is $R$-linear and hence
has a characteristic polynomial
$P(T) = T^n + r_{n-1}T^{n-1} + \ldots + r_0 \in R[T]$.
For any prime
$\mathfrak{p} \in \Spec(R)$, $f$ acts nilpotently on $A
\otimes_R \kappa(\mathfrak{p})$ if and only if $\mathfrak p \in
V(r_0, \ldots, r_{n-1})$.
Moreover, the image of $\Spec(A_f)$ in $\Spec(R)$ is
$\bigcup_{i=0}^{n-1}D(r_i)$. This formula commutes with every ring
base change $R\to R'$, and $\Spec(A)\to\Spec(R)$ is open after every
such base change.
\end{lemma}

\begin{proof}
This follows quite easily once we prove that the characteristic
polynomial $\bar P(T) \in \kappa(\mathfrak p)[T]$ of the
multiplication map $m_{\bar f}: A \otimes_R \kappa(\mathfrak p) \to
A \otimes_R \kappa(\mathfrak p)$ which multiplies elements of $A
\otimes_R \kappa(\mathfrak p)$ by $\bar f$, the image of $f$ viewed in
$A\otimes_R\kappa(\mathfrak p)$, is just the image of $P(T)$ in
$\kappa(\mathfrak p)[T]$. Let $(a_{ij})$ be the matrix of the map
$m_f$ with entries in $R$, using a basis $e_1, \ldots, e_n$
of $A$ as an $R$-module.
Then, $A \otimes_R \kappa(\mathfrak p) \cong (R \otimes_R
\kappa(\mathfrak p))^{\oplus n} = \kappa(\mathfrak p)^n$, which is
an $n$-dimensional vector space over $\kappa(\mathfrak p)$ with
basis $e_1 \otimes 1, \ldots, e_n \otimes 1$. The image $\bar f = f
\otimes 1$, and so the multiplication map $m_{\bar f}$ has matrix
$(a_{ij} \otimes 1)$. Thus, the characteristic polynomial is
precisely the image of $P(T)$.

\medskip\noindent
For a linear endomorphism $N$ of a finite dimensional vector space,
nilpotence implies that the filtration
$0=\ker(N^0)\subset\ker(N)\subset\cdots\subset\ker(N^a)$
eventually reaches the whole space. A basis extending bases along this
filtration gives a strictly triangular matrix, since
$N(\ker(N^j))\subset\ker(N^{j-1})$; its characteristic polynomial
is $T^n$. Conversely, if the characteristic polynomial is $T^n$,
Cayley--Hamilton gives $N^n=0$ for $n>0$.
For $n=0$ the space is zero, its only endomorphism is zero and its
characteristic polynomial is $1=T^0$. Applying this to $m_{\bar f}$
gives nilpotence exactly when all the images of
$r_0,\ldots,r_{n-1}$ in $\kappa(\mathfrak p)$ vanish.
The kernel of $R\to\kappa(\mathfrak p)$ is $\mathfrak p$, proving
the first assertion, including $V()=\Spec(R)$ when $n=0$.

\medskip\noindent
Put $B=A\otimes_R\kappa(\mathfrak p)$. Multiplication by $\bar f$
is nilpotent if and only if $\bar f$ is nilpotent: one implication
follows by applying the power of the multiplication map to $1$, and
the other follows by multiplication by $\bar f^a=0$.
The localization $B_{\bar f}$ is zero if and only if
$\bar f^a=0$ for some $a\geq0$, by the equality criterion for
$1/1=0/1$. Also
$A_f\otimes_R\kappa(\mathfrak p)\cong B_{\bar f}$:
the map sends $(a/f^j)\otimes\lambda$ to
$(a\otimes\lambda)/\bar f^j$, and its inverse is induced by
$a\otimes\lambda\mapsto(a/1)\otimes\lambda$ and the inverse
$(1/f)\otimes1$ of $\bar f$. These maps fix the generators and
the inverse of $f$, so their composites are identities.
Lemma \ref{lemma-in-image} now gives
$$
\mathfrak p\in\operatorname{Im}(\Spec(A_f)\to\Spec(R))
\ \Longleftrightarrow\ B_{\bar f}\ne0
\ \Longleftrightarrow\ \mathfrak p\in\bigcup_{i=0}^{n-1}D(r_i).
$$
For $n=0$, $A=0$ and the union and image are both empty.

\medskip\noindent
For an arbitrary ring map $R\to R'$, the images of the chosen basis
give a basis of $A'=A\otimes_RR'$. Its multiplication matrix for
$f'=f\otimes1$ is the image $(a'_{ij})$ of $(a_{ij})$. The determinant
formula
$$
\det(TI-(a_{ij}))=
\sum_{\sigma\in S_n}\operatorname{sgn}(\sigma)
\prod_{i=1}^n(T\delta_{i,\sigma(i)}-a_{i,\sigma(i)})
$$
shows that every coefficient of its characteristic polynomial is the
image of the corresponding original coefficient, with all signs and
terms retained. Consequently the image of $\Spec(A'_{f'})$ is
$\bigcup_{i=0}^{n-1}D(r'_i)$, the inverse image of
$\bigcup_{i=0}^{n-1}D(r_i)$ in $\Spec(R')$.
Finally, the same image formula applies to every element of $A'$,
whether or not it comes from an element of $A$. Standard opens form
a basis of $\Spec(A')$, and arbitrary unions of their open images
are open. Hence $\Spec(A')\to\Spec(R')$ is open.
\end{proof}